\documentclass[tikz,border=8pt]{standalone}
\usepackage{tikz}
\usepackage{amsmath}
\usepackage{amssymb}
\usepackage{ctex}
\usetikzlibrary{calc, positioning, arrows.meta, decorations.pathreplacing}

\begin{document}
\begin{tikzpicture}[>=Stealth,
  cnode/.style={draw=gray!60, fill=white, rounded corners=5pt,
                minimum width=2.8cm, minimum height=1.1cm,
                inner sep=6pt, align=center, font=\small},
]

%% ===== 四个节点（交换图）=====
% 布局：左下 → 右下 → 左上 → 右上
%   (0,0) = 左下  V_B （输入空间，基 B）
%   (5,0) = 右下  V_{B'}（输入空间，基 B'）
%   (0,4) = 左上  V_B （输出空间，基 B）
%   (5,4) = 右上  V_{B'}（输出空间，基 B'）

\node[cnode] (BLin)  at (0,0) {$V_{\mathcal{B}}$（输入）};
\node[cnode] (BLout) at (0,4) {$V_{\mathcal{B}}$（输出）};
\node[cnode] (BRin)  at (5,0) {$V_{\mathcal{B}'}$（输入）};
\node[cnode] (BRout) at (5,4) {$V_{\mathcal{B}'}$（输出）};

%% ===== 箭头 =====
% 左竖：基 B 下的矩阵 A
\draw[->,very thick,blue!70] (BLin.north) -- (BLout.south)
  node[midway, left, font=\small, blue!80!black] {$A$};

% 右竖：基 B' 下的矩阵 B
\draw[->,very thick,red!70] (BRin.north) -- (BRout.south)
  node[midway, right, font=\small, red!80!black] {$B = P^{-1}AP$};

% 下横：P（输入空间坐标变换）
\draw[->,thick,gray!70] (BLin.east) -- (BRin.west)
  node[midway, below, font=\small, gray!80] {$P$（坐标变换）};

% 上横：P（输出空间坐标变换）
\draw[->,thick,gray!70] (BLout.east) -- (BRout.west)
  node[midway, above, font=\small, gray!80] {$P$};

%% ===== 对角参考线（T 的整体作用）=====
\draw[->,very thick,purple!60,dashed] (BLin.north east)
  to[bend left=10] node[midway, above left, font=\small, purple!70]
  {同一线性变换 $T$} (BRout.south west);

%% ===== 相似矩阵核心公式 =====
\node[fill=purple!8, draw=purple!50, rounded corners=5pt,
      inner sep=8pt, align=center, font=\small]
  at (2.5, 2.0) {
  $B = P^{-1}A\,P$\\[3pt]
  {\footnotesize $A \sim B$（相似）}
};

%% ===== 相似不变量列表（外侧底部）=====
\node[font=\small, fill=white, draw=gray!50, thin, rounded corners=4pt,
      inner sep=10pt, align=center]
  at (2.5, -2.0) {
  \textbf{相似不变量}（$A \sim B$ 时以下量相等）\\[5pt]
  \begin{tabular}{llll}
    迹：& $\operatorname{tr}(A)=\operatorname{tr}(B)$ &
    行列式：& $\det(A)=\det(B)$ \\[3pt]
    特征值：& 相同（含重数） &
    秩：& $\operatorname{rank}(A)=\operatorname{rank}(B)$
  \end{tabular}
};

%% ===== 整体标题 =====
\node[font=\large\bfseries] at (2.5, 5.3) {相似矩阵：同一变换的两种坐标表示};

\end{tikzpicture}
\end{document}
